% Prim's Algorithm

Finds the Minimum Spanning Tree (MST) of a graph from a starting vertex in $O(E \log V)$ using a priority queue. It is optimal for dense graphs and grows a single tree component.

\begin{lstlisting}
const long long INF = 1e18;
struct Edge { int to, weight, id; };

long long prim(int source, int n, const vector<vector<Edge>>& adj, vector<int>& mst_edges) {
    long long total_weight = 0;
    vector<bool> visited(n + 1, false);
    vector<long long> min_edge(n + 1, INF);
    vector<int> parent_edge(n + 1, -1);
    
    // priority_queue stores {weight, vertex}
    priority_queue<pair<long long, int>, vector<pair<long long, int>>, greater<pair<long long, int>>> pq;
    
    min_edge[source] = 0;
    pq.push({0, source});
    int visited_count = 0;

    while (!pq.empty()) {
        auto [w, u] = pq.top(); pq.pop();
        if (visited[u]) continue;
        
        visited[u] = true;
        total_weight += w;
        visited_count++;
        
        if (parent_edge[u] != -1) {
            mst_edges.push_back(parent_edge[u]);
        }

        for (const auto& edge : adj[u]) {
            if (!visited[edge.to] && edge.weight < min_edge[edge.to]) {
                min_edge[edge.to] = edge.weight;
                parent_edge[edge.to] = edge.id;
                pq.push({min_edge[edge.to], edge.to});
            }
        }
    }
    return (visited_count == n) ? total_weight : -1;
}
\end{lstlisting}

To find the MST, build your adjacency list passing an incremental edge ID into \lstinline|Edge.id|. Call \lstinline|prim(1, n, adj, mst_edges)| to compute the minimum total weight and populate the vector with selected edge IDs.
